\chapter{Basic concepts}\label{ch:basics}

\section{Setups, experiments and simulated data}

A GLoBES.jl computation involves
\begin{description}
  \item[A setup] (\code{Setup}): a list of experiments, their simulated data,
    the probability engine, the registered $\chi^2$ functions, and the
    minimiser settings. There is always a current setup, and most programs use
    only that one. Section~\ref{sec:setups} explains how to work with several.
  \item[Experiments] (\code{Experiment}), each loaded from an AEDL file. An
    experiment has one neutrino source, one baseline and matter density
    profile, and one detector with one or more \emph{rules}: independent
    measurements with their own signals, backgrounds and systematics.
  \item[Oscillation parameters] (\code{OscParams}): the point in parameter
    space at which data are simulated or a hypothesis is tested.
\end{description}
The steps of an analysis are always the same:
\begin{julia}
using GLoBES
exp   = Experiment("T2K.glb")                       # 1. load experiments
truth = OscParams(θ12=0.5843, θ13=0.1503, θ23=0.8326,
                  δ=1.5π, Δm21=7.53e-5, Δm31=2.5e-3)
simulate!(truth)                                    # 2. simulate the data
chi2(OscParams(truth; δ=0))                         # 3. test hypotheses
\end{julia}

\section{Oscillation parameters}\label{sec:oscparams}

\begin{table}[t]
\centering\small
\begin{tabular}{lllll}
\toprule
Parameter & Name & ASCII name & Unit & C index \\
\midrule
$\theta_{12}$ & \code{θ12} & \code{theta12} & rad & \code{GLB\_THETA\_12} \\
$\theta_{13}$ & \code{θ13} & \code{theta13} & rad & \code{GLB\_THETA\_13} \\
$\theta_{23}$ & \code{θ23} & \code{theta23} & rad & \code{GLB\_THETA\_23} \\
$\dcp$        & \code{δ}   & \code{delta}   & rad & \code{GLB\_DELTA\_CP} \\
$\dm{21}$     & \code{Δm21} & \code{dm21}   & $\mathrm{eV}^2$ & \code{GLB\_DM\_21} \\
$\dm{31}$     & \code{Δm31} & \code{dm31}   & $\mathrm{eV}^2$ & \code{GLB\_DM\_31} \\
matter density scaling & \code{ρ} & \code{rho} & 1 & density parameters \\
\midrule
$\st{12}$, $\st{13}$, $\st{23}$ & \code{sin²θ12} \dots & \code{sin2theta12} \dots & & \\
$\sst{12}$, $\sst{13}$, $\sst{23}$ & \code{sin²2θ12} \dots & \code{sin22theta12} \dots & & \\
\bottomrule
\end{tabular}
\caption{Names of the oscillation parameters. The last two rows are derived
quantities, which can be read (\code{p.sin²2θ13}) and set with
\code{with\_param} and in scans, but are not stored.}
\label{tab:params}
\end{table}

The standard parameters are listed in Table~\ref{tab:params}. They follow the
standard parametrisation of the leptonic mixing matrix,
\begin{equation}
U = \begin{pmatrix} 1&0&0\\ 0&c_{23}&s_{23}\\ 0&-s_{23}&c_{23}\end{pmatrix}
    \begin{pmatrix} c_{13}&0&s_{13}e^{-i\dcp}\\ 0&1&0\\ -s_{13}e^{i\dcp}&0&c_{13}\end{pmatrix}
    \begin{pmatrix} c_{12}&s_{12}&0\\ -s_{12}&c_{12}&0\\ 0&0&1\end{pmatrix},
\label{eq:pmns}
\end{equation}
with $s_{ij}=\sin\theta_{ij}$, $c_{ij}=\cos\theta_{ij}$, and
$\dm{ij}=m_i^2-m_j^2$. The mass ordering is normal for $\dm{31}>0$ and inverted
for $\dm{31}<0$.

\paragraph{Creating parameters.} All of $\theta_{12}$, $\theta_{13}$,
$\theta_{23}$, $\dm{21}$ and $\dm{31}$ must be given; $\dcp$ defaults to 0 and
$\rho$ to 1:
\begin{julia}
p = OscParams(θ12=asin(sqrt(0.307)), θ13=asin(sqrt(0.022)), θ23=asin(sqrt(0.546)),
              δ=1.5π, Δm21=7.53e-5, Δm31=2.5e-3)
q = OscParams(p; δ=0, Δm31=-2.4e-3)       # a copy with some parameters changed
r = with_param(p, :sin²2θ13, 0.085)        # set θ13 through sin²2θ13
p.θ23, p.sin²θ23, p.ρ                     # read parameters
\end{julia}
\code{OscParams} are immutable: functions never change them, and a changed copy
is made with \code{OscParams(p; changes...)} or \code{with\_param}.

\paragraph{Matter density.} \code{ρ} scales the matter density profile of the
experiments: the densities of the AEDL file are multiplied by $\rho$ when rates
are computed. It is one number for all experiments, or a vector with one entry
per experiment, \code{OscParams(p; ρ=[1.0, 0.97])}. Its uncertainty is taken
into account by fitting it with a prior (Section~\ref{sec:priors}).

\paragraph{Additional parameters.} A probability engine for non-standard
physics (Chapter~\ref{ch:nonstandard}) may have more parameters. They are given
as a named tuple, \code{OscParams(p; extra=(ε=0.1,))}, and accessed by name like
the standard ones.

\section{Units and the integrated luminosity}\label{sec:luminosity}

GLoBES.jl uses GeV for neutrino energies, km for baselines, $\mathrm{g/cm^3}$
for matter densities, $\mathrm{eV}^2$ for mass squared differences and radians
for angles. The number of events is proportional to the \emph{integrated
luminosity},
\begin{equation}
N \propto \text{source power} \times \text{running time} \times \text{target mass},
\end{equation}
with units chosen in the AEDL file: typically MW of beam power (or useful
parent decays per year for neutrino factories and beta beams), years, and kt
(t for reactor detectors). Section~\ref{sec:aedl-flux} explains how the
normalisation of the flux tables relates these units to event numbers. At
runtime, they are changed with \code{source\_power!}, \code{running\_time!} and
\code{target\_mass!} (Section~\ref{sec:exposure}).

\section{Simulated data}\label{sec:simulate}

\begin{apifun}{simulate!(truth::OscParams; setup=current\_setup())}
Compute the event rates of all experiments of the setup for the parameters
\code{truth}. They serve as the data that all $\chi^2$ values compare to.
\end{apifun}

\noindent The simulated data are the event rates expected for the true
parameters; they are not fluctuated. This gives the $\Delta\chi^2$ of an
\emph{average} experiment (the ``Asimov data set''): the expected sensitivity.
Since the true parameters fit the data perfectly, $\chi^2_\text{min}=0$ and
$\chi^2$ and $\Delta\chi^2$ are the same.

\code{simulate!} must be called again after experiments are added or changed;
\code{chi2} and \code{marginalize} throw an error if no data have been
simulated. \code{simulated\_params()} returns the parameters of the last call.

\section{Errors}

The functions of the high-level API throw a \code{GLoBESError} when something is
wrong: a file that cannot be read, an unknown rule name, a missing
\code{simulate!}, a parameter vector of the wrong length. Errors in AEDL files are
\code{AEDLError}s with the file name and line. The \code{glb*} functions instead
print errors and return error codes, like the C library; their output is
controlled with \code{glbSetVerbosityLevel}.

\section{Several setups}\label{sec:setups}

Each \code{Setup} is an independent simulation. Every function works on the
current setup unless it gets a \code{setup} keyword or an experiment, which
belongs to the setup it was loaded into:
\begin{julia}
near = Setup()
far  = Setup()
Experiment("NOvA.glb"; setup=near)
Experiment("NOvA.glb"; setup=far, BASELINE=1500)   # BASELINE: an AEDL variable
simulate!(truth; setup=near); simulate!(truth; setup=far)
chi2(test; setup=near), chi2(test; setup=far)
\end{julia}
\code{with\_setup(f, setup)} makes \code{setup} the current setup while \code{f}
runs, which is convenient for longer pieces of code and also applies to the
\code{glb*} functions:
\begin{julia}
with_setup(far) do
    simulate!(truth)
    glbChiSys(glbparams, GLB_ALL, GLB_ALL)   # the glb* layer sees `far` too
end
\end{julia}
\code{current\_setup()} returns the current setup.

\begin{note}
GLoBES.jl is not thread-safe: functions make their setup the current one while
they run, so two threads computing at the same time interfere even if they use
different setups. To compute in parallel, use several processes (for instance
with the \code{Distributed} standard library), each with its own experiments.
\end{note}

\begin{cdiff}
C GLoBES keeps all its state (experiments, rates, minimiser settings, the
probability engine) in global variables, so a program can only have one
simulation. In GLoBES.jl, this state belongs to a \code{Setup}.
\end{cdiff}
